% Options for packages loaded elsewhere
\PassOptionsToPackage{unicode}{hyperref}
\PassOptionsToPackage{hyphens}{url}
%
\documentclass[
  11pt,
]{article}
\usepackage{amsmath,amssymb}
\usepackage{iftex}
\ifPDFTeX
  \usepackage[T1]{fontenc}
  \usepackage[utf8]{inputenc}
  \usepackage{textcomp} % provide euro and other symbols
\else % if luatex or xetex
  \usepackage{unicode-math} % this also loads fontspec
  \defaultfontfeatures{Scale=MatchLowercase}
  \defaultfontfeatures[\rmfamily]{Ligatures=TeX,Scale=1}
\fi
\usepackage{lmodern}
\ifPDFTeX\else
  % xetex/luatex font selection
\fi
% Use upquote if available, for straight quotes in verbatim environments
\IfFileExists{upquote.sty}{\usepackage{upquote}}{}
\IfFileExists{microtype.sty}{% use microtype if available
  \usepackage[]{microtype}
  \UseMicrotypeSet[protrusion]{basicmath} % disable protrusion for tt fonts
}{}
\makeatletter
\@ifundefined{KOMAClassName}{% if non-KOMA class
  \IfFileExists{parskip.sty}{%
    \usepackage{parskip}
  }{% else
    \setlength{\parindent}{0pt}
    \setlength{\parskip}{6pt plus 2pt minus 1pt}}
}{% if KOMA class
  \KOMAoptions{parskip=half}}
\makeatother
\usepackage{xcolor}
\usepackage[margin=1in]{geometry}
\setlength{\emergencystretch}{3em} % prevent overfull lines
\providecommand{\tightlist}{%
  \setlength{\itemsep}{0pt}\setlength{\parskip}{0pt}}
\setcounter{secnumdepth}{-\maxdimen} % remove section numbering
\ifLuaTeX
  \usepackage{selnolig}  % disable illegal ligatures
\fi
\usepackage{bookmark}
\IfFileExists{xurl.sty}{\usepackage{xurl}}{} % add URL line breaks if available
\urlstyle{same}
\hypersetup{
  pdftitle={All real references on products of rank-one balanced ray stars},
  pdfauthor={P05.2 reviewed partial result},
  hidelinks,
  pdfcreator={LaTeX via pandoc}}

\title{All real references on products of rank-one balanced ray stars}
\author{P05.2 reviewed partial result}
\date{3 October 2026}

\begin{document}
\maketitle

\subsection{Result and exact scope}\label{result-and-exact-scope}

The original arbitrary-balanced-pair classification P05.2 remains open.
This draft settles the specified \textbf{all-real-reference converse on
products of rank-one balanced ray stars times complete affine space}.
The new argument removes the actual ambient-extension assumption from
the previous real-reference draft. It does not assert that such an
extension exists. The mechanism is an exact finite convex decomposition
of a real PA PC reference into rational PA PC references, followed by
addition of their rational-domain modifications. The exact proof and
its complete supporting argument received the source-aware independent
PASS accompanying this paper. The original classification remains open;
no formal verification or priority certification is claimed.

For each \(1\leq l\leq k\), let \(N_l\) be a lattice,
\(V_l=N_l\otimes\mathbb R\), and let \(v_{li}\in N_l\),
\(1\leq i\leq r_l\), be primitive generators of distinct rays spanning
\(V_l\). Assume \[
r_l\geq2,\qquad \dim V_l=r_l-1,\qquad
\ker\bigl((a_i)\mapsto\textstyle\sum_i a_iv_{li}\bigr)
=\mathbb Rw_l,\quad w_l\in\mathbb Q_{>0}^{r_l}.
\tag{S}
\] Set \(G_l=\bigcup_i\mathbb R_{\geq0}v_{li}\) and assign its \(i\)th
ray weight \(\alpha_{li}=c_lw_{li}\), where \(c_l>0\) is real. Let
\(E=N_E\otimes\mathbb R\) have dimension \(m\) and constant weight
\(e>0\). Use the actual direct lattice and carrier \[
N=\bigoplus_{l=1}^kN_l\oplus N_E,\quad W=N\otimes\mathbb R,
\quad X=\prod_{l=1}^kG_l\times E,\quad d=k+m.
\] The top product weights are \(e\prod_l\alpha_{li_l}\). Empty products
are permitted. A factor is \textbf{branched} if \(r_l\geq3\). The
tangent lattice of a cell is always its intersection with \(N\), and
primitive quotient normals are computed in these saturated lattices. The
subgroup collectively generated by a star's rays need not be saturated.

Here \(\mathrm{PA}_{\mathbb Q}\) means continuous functions with
rational affine coefficients on a finite rational spatial complex. The
class \(\mathrm{PA}_{\mathbb R}\) permits real coefficients on a finite
\textbf{rational spatial} complex. The literal source class
\(\mathrm{PA}_{\mathbb Z}\) requires integral tangent slopes and
rational constants. These three conventions are distinct.

\textbf{Theorem (all-real-reference product criterion; partial P05.2).}
Let \(\gamma\in\mathrm{PA}_{\mathbb R}(X)\cap\mathrm{PC}(X)\) be
positive at infinity. The following package holds if and only if at most
one star is branched:

For every bounded \(u\in\mathrm{PA}_{\mathbb R}(X)\), the raw envelope
\(P_\gamma(u)\) is real-regularizable and \[
\int_X\bigl(P_\gamma(u)-(\gamma+u)\bigr)
\mathrm{MA}_{\mathrm{poly},\mathbb R}(P_\gamma(u))=0.
\tag{T}
\] The positive implication is the precise real-reference positive
theorem supplied in caller-imported note \textbf{n100}. Its hypotheses
are checked below; its negative rational-reference quantifier is not
used as an all-real-reference theorem.

The new negative implication is stronger than failure of whole-pair
regularity. If at least two stars are branched, \textbf{every} reference
in the stated real class has a bounded, genuinely nonseparable real-PA
obstacle \(u\), on finite rational spatial domains, whose envelope is an
explicitly constructed \[
P_\gamma(u)=P\in\mathrm{PAPC}_{\mathbb R}(X,\gamma)
\subset\mathrm{PC}_{\mathrm{reg},\mathbb R}(X,\gamma),
\] and whose \textbf{complete} contact defect is \[
\int_X(P-(\gamma+u))\mathrm{MA}_{\mathrm{poly},\mathbb R}(P)
=-\epsilon\lambda^d\kappa<0.
\tag{N}
\] Here \(\epsilon\in\{1,2\}\), \(\lambda>0\) is rational, and the
positive constant computed directly below is \[
\kappa=2^{m+2}d!\,CC'e\prod_{l=3}^k M_l.
\tag{K}
\] No ambient convex extension of \(\gamma\), global PC approximation
theorem, or all-obstacle regularity of a negative pair is assumed.

\subsection{Definitions, exact source inputs, and dependency
boundary}\label{definitions-exact-source-inputs-and-dependency-boundary}

BMPS means Botero--Mazzon--Pille-Schneider, \emph{Monge--Ampère measures
on balanced polyhedral spaces}, arXiv:2603.08664v1, submitted 9 March
2026, manuscript dated 10 March 2026. The cached PDF is snapshotted as
\texttt{sources/BMPS2603.08664v1.pdf}; its SHA-256 is
\texttt{5cf7e68ee9113f1dc813f9aefa5f9ca72cba02fc4fa688f1a61e4b7f0d1223e7}.
The extracted body and the inspected numbered passages are retained
locally.

Definition 2.1 fixes saturated tangent lattices; Definitions 2.18--2.20 use
\(N_\sigma=N\cap\operatorname{span}_{\mathbb R}(\sigma-\sigma)\) and
normal lifts \(n_{\sigma/\tau}\) whose quotient classes primitively
point from \(\tau\) into \(\sigma\). Definitions 2.19 and 2.20 use
\textbf{rational} weights and cycles. We explicitly extend weights to
real scalars; balancing means
\(\sum_{\sigma\succ\tau}z(\sigma)n_{\sigma/\tau}\in N_{\tau,\mathbb R}\).
Definition 2.22 permits generalized fans with lineality. Definition 2.28
is the literal integral-slope/rational-constant PA class, not our real
category.

Definitions 3.2--3.3 give the actual PCC rule: endpoints
\(x_j\in X\cap N_{\mathbb Q}\), coefficients \(a_j\in[0,1]\) real,
\(\sum a_j=1\), centre \(x=\sum a_jx_j\), and cells
\(\tau\preceq\sigma_j\) in some rational structure with \(x\in\tau\),
\(x_j\in\sigma_j\). A finite raw PC function satisfies
\(h(x)\leq\sum a_jh(x_j)\) for all such combinations. No continuity
theorem for raw competitors is used. Every PA function in this paper is
continuous by definition.

Definition 3.1 tests PAPC against \textbf{all} nonnegative local
balanced cycles of every positive dimension on ordinary open sets. We
prove the needed PA-PC effectivity in the real category below. We do not
simply extend the source's Lemma 3.4 without checking its coefficient
hypotheses.

Definition 2.31 gives the divisor formula, with \(a_\sigma\) denoting an
affine piece's linear part, \[
(a\cdot z)(\tau)=\sum_{\sigma\succ\tau}z(\sigma)a_\sigma(n_{\sigma/\tau})
-a_\tau\left(\sum_{\sigma\succ\tau}z(\sigma)n_{\sigma/\tau}\right).
\tag{D}
\] Remark 2.33 and Lemmas 2.34--2.35 state normal-lift independence,
refinement compatibility, linearity, balancing, and commutativity. Their
inspected proof bodies use cancellation and reduction to two-dimensional
quotient stars, the latter citing Allermann--Rau Proposition 3.7. We
reconstruct the requisite two-dimensional computation and refinement
argument below, so no unretrieved external theorem body is an essential
premise of the new negative proof.

Definition 4.4 identifies the PA operator with \(a^d\cdot[X]\), the
finite atomic zero-cycle. We retain this intersection normalization and
all multinomial factorials. Definition 4.15 requires strictly positive
reference slope on every nonzero recession direction of every affine
reference cell. Definition 3.14 gives the bounded-reference growth
class. We use the real analogue of Proposition 3.22 and Definition
3.23, with real regularizability explicitly defined here by uniform
decreasing real PAPC approximation. In
particular \(P+1/j\) suffices when \(P\) is already PAPC; its measures
are identical.

The envelope of Definition 6.2 is exactly \[
P_\gamma(u)(x)=\sup\{h(x):h\in\mathrm{PC}(X),\ h\leq\gamma+u,
\ h\leq\gamma+A_h\text{ for some }A_h\in\mathbb R\}.
\tag{E}
\] Definition 6.9 assumes regularity before asking for all-obstacle
orthogonality. Our negative witness has a measure in its own right and
disproves the contact identity regardless of other envelopes. BMPS's
uniform operator-extension theorem, total-mass theorem, Alexandrov
comparison, analytic realization and projection formula are not premises
of the negative proof.

The only imported supporting result not reproved here is the
\textbf{positive side} of n100: on the above product class with at most
one branched factor, every positive-at-infinity real-PA PC reference and
every bounded real-PA obstacle have a real-PAPC envelope and zero
contact defect. It includes arbitrary nonseparable obstacles and uses
finite rational spatial domains. Bivalent factors merge with the affine
factor, and the dimension-zero case is included. The supplied complete
argument is \texttt{sources/n100-argument.md}, SHA-256
\texttt{fed51d88551d295be9bbf798a418762699a98b49b7d7747ce4a8517ea2d8b2a3};
its exact-scope review receipt is retained. This is the caller-imported
theorem n100, not a new local note. Its positive reference quantifier is
already real. The previous real-reference draft, self-audit and research
notes were read as fallible material and are retained under
\texttt{sources/prior-*}; none is a supporting theorem. The present
exact-scope review reconstructed the full unchanged positive supporting
proof embedded in its argument, together with the new converse. Both
arguments and the complete review receipt accompany this paper.

\subsection{Elementary finite-dimensional
facts}\label{elementary-finite-dimensional-facts}

The following proofs make the new finite support and decomposition
arguments independent of an unstated polyhedral extension or
approximation theorem.

\textbf{Lemma 1 (finite generation and closedness).} A rational
polyhedral cone, including one with lineality, has finitely many
rational generators. A nonempty pointed rational polyhedron has the form
\(\operatorname{conv}(A)+\operatorname{cone}(R)\) for finite nonempty
rational \(A\) and finite rational \(R\). A finitely generated real cone
is closed. Projection of a finite system of nonstrict real linear
inequalities is a finite system of nonstrict linear inequalities;
rational data remain rational.

\textbf{Proof.} Write a pointed cone as \(\ell_i(z)\geq0\), including
equality constraints as pairs. Its defining forms have zero common
kernel, since a vector in that kernel would give a line. Therefore
\(\ell=\sum_i\ell_i\) is positive on every nonzero cone vector. The
section \(\ell=1\) is bounded: an unbounded sequence, normalized and
passed to a convergent subsequence on the unit sphere, would yield a
nonzero cone vector with \(\ell=0\). The section is a compact
polyhedron. A point in a positive-dimensional minimal face is on a line
in that face with two terminal points in smaller faces; induction
expresses it as a convex combination of vertices. There are finitely
many vertices, each obtained from a full-rank selection of active
equations in the affine hull, and rational equations give rational
vertices. Their directions generate the cone. For a cone with lineality,
quotient by its rational lineality, apply this argument to the pointed
quotient, choose rational lifts, and append both signs of a rational
lineality basis. This yields finite rational generators.

Homogenize a nonempty pointed polyhedron \(P=\{x:Ax\geq b\}\) to
\(\{(y,t):Ay\geq bt,\ t\geq0\}\). A line in this cone would have \(t=0\)
and give lineality of \(P\). Apply the pointed cone result. Generators
with \(t>0\), after dividing by \(t\), give \(A\subset P\); those with
\(t=0\) give \(R\subset\operatorname{rec}(P)\). Taking the section
\(t=1\) proves the assertion, with \(A\) nonempty.

For projection, eliminate a variable by requiring each of its finitely
many lower bounds to be at most each upper bound, retaining inequalities
independent of that variable. A one-sided finite list of bounds is
always satisfiable. This is exact over the reals, preserves nonstrict
inequalities, and preserves rational coefficients. Repetition proves the
projection statement. A cone generated by \(g_i\) is the projection of
\(x=\sum_i t_ig_i\), \(t_i\geq0\), so is a closed polyhedron by that
statement. The argument applies to real generators as well. \(\square\)

\textbf{Lemma 2 (finite support alternative).} Given finite vectors
\(d_i\) and real numbers \(b_i\), there exists a linear form \(p\) with
\(p(d_i)\leq b_i\) for all \(i\) if \[
\sum_i t_ib_i\geq0\quad\text{whenever }t_i\geq0,
\quad\sum_i t_id_i=0.
\tag{F}
\] The converse is immediate.

\textbf{Proof.} Let \(K=\operatorname{cone}\{(d_i,b_i),(0,1)\}\), a
closed cone by Lemma 1. Condition (F) excludes \(z=(0,-1)\) from \(K\).
Let \(k\) be a nearest point of \(K\) to \(z\); existence follows by
restricting to a sufficiently large compact ball. With \(r=k-z\), the
nearest-point inequality gives \(\langle r,w-k\rangle\geq0\) for all
\(w\in K\). Taking \(w=0,2k\) gives \(\langle r,k\rangle=0\). Thus
\(\langle r,w\rangle\geq0\) and \(\langle r,z\rangle=-\|r\|^2<0\). Write
\(r=(-p_0,\alpha)\). The last inequality gives \(\alpha>0\), and each
generator gives \(p_0(d_i)\leq\alpha b_i\). Set \(p=p_0/\alpha\).
\(\square\)

\subsection{A finite rational system for PA polyhedral
convexity}\label{a-finite-rational-system-for-pa-polyhedral-convexity}

Let \(\Omega\) be \textbf{any} fixed finite rational polyhedral complex
in a lattice vector space \(W\), with carrier \(Y\). For a face \(\tau\)
and an incident cell \(\sigma\succeq\tau\), put \[
T_{\sigma,\tau}=\{v:x+tv\in\sigma\text{ for all sufficiently small }t\geq0\},
\qquad x\in\operatorname{RelInt}(\tau).
\] This does not depend on the chosen \(x\): it is given by the
constraints of \(\sigma\) active on \(\tau\). It is a rational
polyhedral cone, with finite rational generators \(d_{\sigma,\tau,q}\)
by Lemma 1. Both signs of the tangent directions of \(\tau\) are
included through the cone's lineality. If \(a\) is PA on \(\Omega\),
write \(a_{\sigma,\mathrm{lin}}\) for its derivative on the tangent
space of \(\sigma\).

\textbf{Lemma 3 (local support characterizes actual PCC convexity).} A
continuous real PA function \(a\) on \(\Omega\) is PC on \(Y\) if and
only if, for every \(\tau\in\Omega\), there exists an ambient linear
form \(p_\tau\in W^*\) with \[
p_\tau(d_{\sigma,\tau,q})
\leq a_{\sigma,\mathrm{lin}}(d_{\sigma,\tau,q})
\quad\text{for all }\sigma\succeq\tau,q.
\tag{L}
\] These are finitely many rational linear inequalities in the affine
coefficients of \(a\) and the coefficients of \(p_\tau\), together with
rational linear continuity equations for \(a\).

\textbf{Proof, necessity.} Pick
\(x\in\operatorname{RelInt}(\tau)\cap N_{\mathbb Q}\). A common rational
\(t>0\) makes every \(x+td_i\) rational and in its incident cell, where
\(i\) runs through the finite generators at \(\tau\). For any
\(t_i\geq0\) with \(\sum t_id_i=0\) and \(S=\sum t_i>0\), these points
and coefficients \(t_i/S\) have centre \(x\). The face \(\tau\) is
incident to every endpoint cell, so this is an actual PCC on \(\Omega\).
Affinity gives \(0\leq\sum_i t_i a_{\sigma_i,\mathrm{lin}}(d_i)\). The
case \(S=0\) is trivial. Lemma 2 supplies \(p_\tau\). The argument
explicitly allows real coefficients in this PCC, exactly as Definition
3.2 does, and uses only rational endpoints. There is no intermediate
centre off the carrier.

\textbf{Proof, sufficiency.} At every
\(x\in\operatorname{RelInt}(\tau)\), (L) extends by nonnegative linear
combination to all of \(T_{\sigma,\tau}\). For every incident affine
piece and sufficiently small feasible \(v\), it gives \[
a(x+v)\geq a(x)+p_\tau(v).
\tag{G}
\] All cells containing sufficiently near points contain \(\tau\): cells
not containing \(x\) are a positive distance from \(x\), there are
finitely many cells, and a cell containing a relative-interior point of
\(\tau\) contains \(\tau\). Thus (G) is a simultaneous local support
inequality on the carrier germ.

Consider any line segment \([x,y]\subset Y\). Restriction of \(a\) is
continuous PL with finitely many intervals. At each interior breakpoint
\(z\), use its minimal face and its form \(p\). If \(v=y-x\), the left
derivative \(b_-\) and right derivative \(b_+\) satisfy
\(b_-\leq p(v)\leq b_+\) by (G) in the two directions. Hence slopes
along the segment are nondecreasing and the restriction is convex. At
its initial point, the first directional derivative is at least
\(p_\tau(y-x)\). Integration of the nondecreasing slopes yields \[
a(y)\geq a(x)+p_\tau(y-x).
\tag{S1}
\] Now take \textbf{any} actual PCC on \textbf{any} rational structure
of \(Y\). Each endpoint is in a convex cell containing the centre, so
every segment from the centre to an endpoint lies in \(Y\). Apply (S1)
with the single local form at the centre, multiply by the PCC
coefficients and sum. The linear terms sum to zero. This proves exactly
the PCC inequality, including arbitrary real centres and coefficients,
without a routing assumption.

For rationality of the system, choose rational basepoints and rational
tangent bases in the affine hulls of the cells. Intrinsic affine
assignments are finite vectors of real coefficients. Agreement on each
shared face is given by rational linear equations (evaluate at a
rational basepoint and rational tangent basis). Each derivative at a
rational \(d_i\) is a rational-linear expression in these variables. The
\(p_\tau(d_i)\) expressions are also rational-linear in the ambient
coefficients of \(p_\tau\). This proves the asserted finite system.
\(\square\)

The lemma is a statement about \textbf{PA functions on a fixed finite
rational complex}. It makes no approximation or regularization claim for
non-PA PC functions or arbitrary envelopes. The local forms need not be
global ambient supporting forms for \(a\).

\textbf{Lemma 4 (real PC implies real PAPC at the scope used here).} A
continuous real PA PC function on finite rational domains has
nonnegative divisor against every nonnegative real local balanced cycle
of every positive dimension.

\textbf{Proof.} Take a common rational refinement of the function and a
local cycle. At a tested codimension-one face \(\tau\) and a
relative-interior point, Lemma 3 gives a simultaneous local support form
\(p\). Both tangent directions along \(\tau\) show
\(p=a_{\tau,\mathrm{lin}}\) on its tangent space. Every inward normal
derivative of \(a_\sigma-p\) is nonnegative by (G); tangential changes
of a normal do not change this conclusion. Formula (D), with the ambient
linear form subtracted, becomes \[
(a\cdot z)(\tau)=\sum_{\sigma\succ\tau}z(\sigma)
(a_{\sigma,\mathrm{lin}}-p)(n_{\sigma/\tau})\geq0.
\] This applies on ordinary open sets as well, because each tested germ
is unchanged by restriction to the open set. It tests all dimensions,
all local cycles and real weights. \(\square\)

\subsection{Exact rational convex decomposition of a real
reference}\label{exact-rational-convex-decomposition-of-a-real-reference}

\textbf{Lemma 5 (exact decomposition, preserving strict growth).} If
\(\gamma\in\mathrm{PA}_{\mathbb R}(Y)\cap\mathrm{PC}(Y)\) on a fixed
finite rational complex is positive at infinity, then there are finitely
many \(\gamma_j\in\mathrm{PA}_{\mathbb Q}(Y)\cap\mathrm{PC}(Y)\), each
positive at infinity on that same complex, and real \(a_j>0\) with
\(\sum_j a_j=1\), such that \[
\gamma=\sum_{j=1}^s a_j\gamma_j\quad\text{exactly on all }Y.
\tag{R}
\] There is no approximation error, including at infinity.

\textbf{Proof.} By Lemma 3, choose a feasible lifted vector \(y\)
consisting of all affine coefficients of \(\gamma\) and all support
forms \(p_\tau\). It satisfies finitely many rational homogeneous
equations \(Ey=0\) and nonstrict inequalities \(Ay\geq0\).

Every affine reference cell is pointed: nonzero lineality would give two
opposite recession directions with opposite slopes, contradicting strict
positivity. By Lemma 1, each recession cone has finitely many rational
generators. Strict positivity at infinity is therefore exactly a finite
list of strict rational-linear inequalities \(By>0\), one for every
nonzero generator. Indeed positive linear values on these generators are
positive on every nonzero nonnegative combination. Bounded cells add no
condition.

Add to \(Ey=0\) the equalities given by precisely those rows of \(A\)
that vanish at \(y\), and call the resulting kernel \(L\). It is a
rational linear subspace. Within \(L\), all remaining \(A\) inequalities
and all \(B\) inequalities hold strictly near \(y\). Choose a rational
basis of \(L\) and a sufficiently small coordinate box about \(y\), with
rational endpoints strictly enclosing each coordinate of \(y\). Its
entire box satisfies those strict inequalities. All active \(A\) rows
remain zero. Every vertex \(y_j\) is rational, satisfies the lifted
support and continuity equations, and has strict growth. Standard box
barycentric weights are products of the one-dimensional weights \$
(b\_i-y\_i)/(b\_i-a\_i)\$ and \((y_i-a_i)/(b_i-a_i)\); they are positive
real numbers summing to one and reproduce every coordinate of \(y\).
Thus \(y=\sum_j a_jy_j\) exactly. If \(L\) has dimension zero, \(y=0\)
is rational and the singleton representation suffices whenever the
strict conditions are empty; with nonempty strict growth this case is
impossible.

Project each \(y_j\) to its affine coefficients \(\gamma_j\). These give
rational affine functions on the same spatial complex. Lemma 3 proves
their actual PC inequalities; the finite strict generator tests prove
their growth. Projection of the exact vector identity gives (R).
\(\square\)

Although separate \(\gamma_j\) can have different recession slopes,
their \textbf{exact weighted sum} is the original reference. The
construction below adds exact modifications of all components and
therefore restores the prescribed exterior function identically. Merely
approximating \(\gamma\) by a rational reference would not do this.

\subsection{Recursive divisor algebra and
normalizations}\label{recursive-divisor-algebra-and-normalizations}

We give the algebra needed for the explicit mass computation. Formula
(D) is bilinear. Normal-lift changes are tangent to \(\tau\) and cancel
by continuity. For an ambient affine function, its linear part pulls out
of the sum in (D), so its divisor on a balanced cycle is zero. The
formula depends only on incident germs, which proves recursive locality.

Refinement is checked locally. If a new tested face lies in the interior
of an old supported cell, the two adjacent subdivisions have equal
inherited weights, equal affine pieces and opposite primitive normals,
so contribute zero. If its minimal old face has the same dimension as
the tested face, the saturated tangents, normal quotient classes and
inherited incident contributions agree with the old formula. If that
minimal face has dimension larger than the original cycle's dimension,
all inherited weights vanish. These exhaust the possibilities and prove
refinement compatibility, including local cycles and generalized fans.

For balancing and commutativity, quotient a codimension-two face by its
saturated tangent lattice, subtract the functions' common ambient linear
extensions along that tangent, and refine to pointed two-cones if
necessary. For each two-cone with boundary primitive rays \(u_r,u_s\),
write its weight \(\omega_{rs}\) and its saturated plane index \[
\chi_{rs}=[N_{C_{rs}}:\mathbb Zu_r+\mathbb Zu_s].
\] The inward primitive normal at \(r\) is \(u_s/\chi_{rs}\) modulo
\(\mathbb Ru_r\). Balancing at \(r\) has the form \[
\sum_{C_{rs}\succ r}\frac{\omega_{rs}}{\chi_{rs}}u_s=t_ru_r.
\] Tangential terms in normal lifts cancel in (D), leaving first divisor
weight \[
D_a(r)=\sum_{C_{rs}\succ r}\frac{\omega_{rs}}{\chi_{rs}}a(u_s)-t_ra(u_r).
\] Summing \(D_a(r)u_r\) and substituting the balancing expression for
\(t_ru_r\) cancels, for each unordered cone, the two terms
\(a(u_s)u_r+a(u_r)u_s\). Thus the first divisor is balanced. A second
divisor at the quotient origin is \[
\sum_{\{r,s\}}\frac{\omega_{rs}}{\chi_{rs}}
\bigl(a(u_s)b(u_r)+a(u_r)b(u_s)\bigr)
-\sum_rt_ra(u_r)b(u_r),
\] which is symmetric in \(a,b\). This proves commutativity at every
tested codimension-two face. Refinement allows the computation even when
the initial germ has lineality. Saturation is retained in the quotient,
so \(\chi_{rs}\) is the actual normal index. These arguments work
directly for real coefficients and weights; no cycle pushforward is
used.

For cycles \(Z,Y\) in actual direct lattice summands, product tangent
lattices split. A codimension-one product face loses a dimension in
exactly one factor, with normal \((n,0)\) or \((0,n)\). For a function
pulled back from the first factor, (D) gives \((a\cdot Z)\times Y\) on
the first type and zero on the second type. Hence \[
(a\circ p_Z)\cdot(Z\times Y)=(a\cdot Z)\times Y.
\tag{PR}
\] This includes zero cycles, lineality, refinements and every
successive intersection. The direct ambient lattice splitting prevents
any extra product index.

\subsection{The two-branched-star trap and its full
mass}\label{the-two-branched-star-trap-and-its-full-mass}

Relabel two branched stars as \(G_1,G_2\). Write their relations
\(w_i,w'_j\), weights \(\alpha_i=c_1w_i\), \(\beta_j=c_2w'_j\).
Partition each set of rays into three nonempty groups. Put
\(W_g=\sum_{i\in g}w_i\), \(W'_g=\sum_{j\in g}w'_j\), and \[
u_1=(1,0),\quad u_2=(0,1),\quad u_3=(-1,-1).
\] The assignment \(L_0(v_i)=u_g/W_g\) annihilates the unique relation
and defines a rational linear map. Likewise define \(L'_0\). Choose one
positive integer \(D\) clearing both matrices in lattice bases, and set
\(L=DL_0\), \(L'=DL'_0\). Then \[
L(v_i)=k_iu_g,\quad k_i=D/W_g\in\mathbb Z_{>0},\qquad
L'(v'_j)=\ell_ju_g,\quad\ell_j=D/W'_g\in\mathbb Z_{>0},
\] where scalar integrality follows from integrality of the maps and
primitivity of \(u_g\). Each group satisfies \[
\sum_{i\in g}\alpha_ik_i=C=c_1D,\qquad
\sum_{j\in g}\beta_j\ell_j=C'=c_2D.
\tag{G1}
\]

On \(G_1\times G_2\), set \(F_1=(La)_1-(L'b)_1\),
\(F_2=(La)_2-(L'b)_2\), \(A=|F_1|\), \(B=|F_2|\), \(q=A+B\). These
restrict ambient rational convex homogeneous functions. In quadrant
\((a,b)=(tv_i,sv'_j)\), \(F_1=p_it-p'_js\), where \(p_i=k_i,0,-k_i\)
according to groups 1,2,3. Interior bends of \(A\) occur precisely when
both rays belong to group 1 or both to group 3. Put
\(g_{ij}=\gcd(k_i,\ell_j)\) and \[
h_{ij}=(\ell_jv_i/g_{ij},\ k_iv'_j/g_{ij}).
\] The saturated quadrant tangent lattice is
\(\mathbb Z(v_i,0)\oplus\mathbb Z(0,v'_j)\). The displayed diagonal
vector is primitive there and in the ambient direct lattice. The
primitive quotient covector is \((k_i/g_{ij},-\ell_j/g_{ij})\). Bezout
supplies a primitive inward normal lift on which \(k_it-\ell_js\) has
value \(g_{ij}\); opposite normals and the absolute-value bend give
weight \(2g_{ij}\alpha_i\beta_j\).

On a nonzero second-factor axis in group 2 the outgoing normal
derivatives give \(\beta_j\sum_i\alpha_i|p_i|=2C\beta_j\). On group 1 or
3 axes, the sign is fixed and the derivative sum is
\(\pm\beta_j\sum_i\alpha_ip_i=0\). First-factor axes give
\(2C'\alpha_i\) in group 2 and zero otherwise. Other quadrant interiors
are affine. This includes every first-divisor face: \[
\begin{split}
A\cdot[G_1\times G_2]={}&
\sum_{g=1,3}\sum_{i,j\in g}2g_{ij}\alpha_i\beta_j[h_{ij}]\\
&+\sum_{i\in2}2C'\alpha_i[(v_i,0)]
+\sum_{j\in2}2C\beta_j[(0,v'_j)].
\end{split}
\tag{D1}
\] Here \([h]\) is the ray of primitive generator \(h\). The first
projections of its diagonal weighted vectors sum to
\(2C'\sum_{i\in1,3}\alpha_iv_i\), and the axis terms complete this to
zero; the second projections similarly complete
\(2C\sum_{j\in1,3}\beta_jv'_j\) to zero. Thus balancing is also checked
directly.

The function \(A\) vanishes on its entire divisor. The slopes of \(B\)
on its diagonals are zero, on its first axes \(k_i\), and on its second
axes \(\ell_j\). Therefore \[
A^2\cdot[G_1\times G_2]=0,\qquad
B\cdot A\cdot[G_1\times G_2]=4CC'\delta_0.
\] Coordinate exchange gives \(B^2=0\). Bilinearity and the two mixed
orders give \[
q^2\cdot[G_1\times G_2]=8CC'\delta_0.
\tag{D2}
\] This is a direct divisor proof, not a pushforward by the nonproper
difference map.

There is an actual positive balanced \textbf{zero curve} for \(q\): take
all \(h_{ij}\) with both rays in the same group, now including group 2,
and ray weights \[
\theta_{ij}=\alpha_i\beta_jg_{ij}.
\tag{Z}
\] Its first projection sums to \(C'\sum_i\alpha_iv_i=0\) and its second
to \(C\sum_j\beta_jv'_j=0\). Thus \(\sum\theta_{ij}h_{ij}=0\). Its
weights have the common real factor \(c_1c_2\) times rational numbers.
The function \(q\) is zero on all its rays. The curve will only
manufacture an actual PCC.

For every other star, put \(h_l(tv_{li})=t\). An ambient rational convex
extension is the maximum of forms \(p_{li}\) determined by \[
p_{li}(v_{li})=1,\qquad
p_{li}(v_{lj})=-w_{li}/\sum_{a\ne i}w_{la}\quad(j\ne i).
\] These assignments annihilate the unique relation, define rational
forms, and their maximum restricts to \(h_l\). Primitive outgoing slopes
one give \(h_l\cdot[G_l]=M_l\delta_0\), where
\(M_l=\sum_i\alpha_{li}>0\). A second intersection vanishes. Choose a
lattice basis of \(N_E\) with dual coordinate functions \(z_a\). On each
corresponding affine coordinate line, \(|z_a|\) has first divisor
\(2\delta_0\) and second divisor zero.

On the full product define \[
Q=q+\sum_{l=3}^kh_l+\sum_{a=1}^m|z_a|,
\qquad \rho=\sum_{l=1}^kh_l+\sum_{a=1}^m|z_a|.
\tag{Q}
\] Both are rational homogeneous continuous PA PC functions. The
function \(\rho\) is proper, positive off zero, and comparable with an
ambient norm on the finite product cones. There is a rational \(K>0\)
with \(0\leq Q\leq K\rho\), by taking finite rational-generator slope
bounds. Appending zero coordinates, the curve (Z) remains a zero curve
for \(Q\).

By (PR), the only nonzero exponent vector in the expansion of \(Q^d\) is
\((2,1,\ldots,1)\): two \(q\) factors, one factor for each remaining
star and each affine coordinate. Any other exponent vector summing to
\(d=2+(k-2)+m\) exceeds one factor's dimension. The multinomial
multiplicity is \(d!/2!\). Thus the \textbf{full} recursive measure is
\[
Q^d\cdot[X]=\frac{d!}{2}\,8CC'\,e\,2^m\prod_{l=3}^kM_l\,\delta_0
=\kappa\delta_0.
\tag{M}
\] All axis, diagonal, branching, lower-face and mixed-order terms are
included. This retains the intersection factorial \(d!\); no probability
normalization or top-face-interior Alexandrov substitution is made.

\subsection{Gluing the trap to every real reference with exact exterior
growth}\label{gluing-the-trap-to-every-real-reference-with-exact-exterior-growth}

Fix an arbitrary reference \(\gamma\) in the negative theorem. Apply
Lemma 5 on any finite rational structure for it: \[
\gamma=\sum_{j=1}^s a_j\gamma_j,\qquad a_j>0,\quad\sum a_j=1,
\] where every \(\gamma_j\) is rational PA PC and positive at infinity.
This exact identity holds also on all unbounded cells.

Refine their structure by the rational domains of \(\rho\). Cells remain
pointed. Write them as finite rational convex hulls plus cones by Lemma
1. On every nonzero recession generator, the \(\rho\) slope is positive
and each \(\gamma_j\) slope is positive. Choose a common rational
\(\eta>0\) below the finite list of ratios. Then \(\gamma_j-\eta\rho\)
has nonnegative recession slopes and is bounded below by its finitely
many point-generator values. For a common rational \(B>0\), \[
\gamma_j\geq\eta\rho-B\quad\text{on }X\quad(1\leq j\leq s).
\tag{C}
\] Bounded cells contribute only point generators. Positivity of the
\(\rho\) recession slope follows directly from its homogeneous radial
formula on each refined cone.

Choose rational \(R>0\) and rational \(\lambda>0\) with
\(\lambda K<\eta\). For each \(j\), choose a rational \(M_j>0\) so large
that \(\gamma_j-M_j<\lambda Q\) on the compact set \(\rho\leq2R\).
Define \[
P_j=\max\{\lambda Q,\gamma_j-M_j\},\qquad
P=\sum_j a_jP_j,\qquad \overline M=\sum_ja_jM_j.
\tag{A}
\] Finite maxima preserve actual PCC inequalities: choose a maximizing
function at the centre and bound it by the endpoint maxima. Therefore
every \(P_j\) and \(P\) is PC. Each \(P_j\) is rational PA on a finite
rational subdivision because all its comparison walls are rational. The
finite common refinement of these subdivisions makes \(P\) real PA;
addition introduces no new comparison wall.

Every \(P_j-\gamma_j=\max(\lambda Q-\gamma_j,-M_j)\) is bounded. Hence
\(P-\gamma\) is bounded exactly, by the original decomposition. More
precisely (C) yields \[
\gamma_j-M_j-\lambda Q\geq(\eta-\lambda K)\rho-B-M_j.
\] For a common rational \(S>2R\) this is nonnegative on \(\rho\geq S\),
so \[
P=\lambda Q\quad(\rho\leq2R),\qquad
P=\gamma-\overline M\quad(\rho\geq S).
\tag{EX}
\] This is exact exterior agreement with the prescribed real reference
up to one constant, not asymptotic slope approximation. Lemma 4 proves
\(P\) is real PAPC against all local cycles. The sequence \(P+1/j\)
decreases uniformly to it in the same growth class, with identical
divisors, establishing the actual regularizable domain.

This gluing uses no ambient convex extension and no real-coefficient
maximum with \(\gamma\). In particular \(\max(\lambda Q,\gamma-M)\) is
\textbf{not} declared PA on rational domains. It may have irrational
walls. Only rational-component maxima are used.

\subsection{The raw envelope, nonseparable obstacle, and complete
negative
integral}\label{the-raw-envelope-nonseparable-obstacle-and-complete-negative-integral}

Put \[
b_R=(1-\rho/R)_+,\qquad f=P+\epsilon b_R,\qquad u=f-\gamma.
\tag{O}
\] The function \(u\) is bounded continuous real PA on finite rational
domains. The roof is rational PA and compactly supported. The positive
rational integer \(\epsilon\) will be selected as 1 or 2 below.

For each zero-curve generator, append zeros in the remaining factors and
\(E\), and set \[
x_{ij}=\frac R{\rho(h_{ij})}h_{ij},\qquad
b_{ij}=\frac{\theta_{ij}\rho(h_{ij})}{\sum_{i',j'}\theta_{i'j'}\rho(h_{i'j'})}.
\tag{PCC0}
\] These are rational endpoints with \(\rho=R\) and \(Q=0\). The
coefficients are positive rational numbers, because their weights have a
common real scalar times rational numbers. They sum to one, and
\(\sum b_{ij}x_{ij}=0\) by (Z). Each endpoint is in a closed original
product cone incident to the actual common face \(\{0\}^k\times E\)
containing zero. This is a valid PCC even if the affine factor has not
been subdivided at zero. At its endpoints, (EX) gives
\(f=P=\lambda Q=0\). Therefore every raw competitor \(h\leq f\) in (E)
satisfies \(h(0)\leq0=P(0)\).

For rational \(x\) with \(0<\rho(x)<R\), the point \(y=Rx/\rho(x)\) is
rational, belongs to the same closed original product cone, and has
\(\rho(y)=R\), \(f(y)=P(y)=\lambda Q(y)\). The endpoints \(0,y\) and
centre \(x\) form a same-cell PCC. Homogeneity and the origin bound give
\[
h(x)\leq\frac{\rho(x)}R h(y)+\left(1-\frac{\rho(x)}R\right)h(0)
\leq\lambda Q(x)=P(x).
\tag{PCC1}
\] For a real \(x\) in this ball, take its minimal face in a common
rational affine-domain structure for \(Q,\rho\). A zero-dimensional face
is rational. In a positive-dimensional face, a sufficiently small
simplex with rational vertices in that face contains \(x\) and is still
inside \(\rho<R\); construct it by approximating vertices of a small
enclosing simplex in the face's rational affine hull. A same-cell PCC
with its real barycentric coefficients and the rational-point bounds
(PCC1) proves the same inequality exactly by affinity of \(Q\). This
avoids any continuity assumption on raw \(h\). For \(\rho\geq R\),
\(b_R=0\) and \(h\leq f=P\) directly. Conversely \(P\) is a PC
competitor in (E), with bounded reference difference. Consequently \[
P_\gamma(u)=P\quad\text{on all }X.
\tag{EP}
\]

Locality of the recursive divisor and (EX),(M) give \[
\mathbf1_{\{\rho<R\}}\mathrm{MA}_{\mathrm{poly},\mathbb R}(P)
=\lambda^d\kappa\delta_0.
\] The whole measure is positive finite atomic by Lemma 4 and the
divisor algebra. The integrand \(P-f=-\epsilon b_R\) vanishes everywhere
outside this open ball, \textbf{including its boundary}. At zero it is
\(-\epsilon\). Thus the complete integral is exactly (N), including
every atom in the interpolation region and on every lower face outside
the roof ball: all those atoms have zero integrand. Only the raw measure
on \(X\) is used; no compactification boundary measure is introduced. No
assertion of total mass alone replaces this contact computation.

To ensure nonseparability, fix rays in the first two factors and form
the rectangle at \(0\), \((Rv_i,0)\), \((0,Rv'_j)\), \((Rv_i,Rv'_j)\),
with the other coordinates zero. Define \[
\Delta a=a(0)+a(Rv_i,Rv'_j)-a(Rv_i,0)-a(0,Rv'_j).
\] Every sum of functions of separate factors has zero mixed difference,
whereas \(\Delta b_R=1\). Set \(A_0=\Delta(P-\gamma)\) and choose
\(\epsilon=1\) unless \(A_0=-1\), in which case choose \(\epsilon=2\).
Then \(\Delta u=A_0+\epsilon\ne0\). These choices leave all PCC endpoint
values and the measure computation unchanged. The negative theorem is
proved for every real reference in its advertised scope.

\subsection{Positive side and completion of the product
criterion}\label{positive-side-and-completion-of-the-product-criterion}

Our product hypotheses are precisely those of the positive result
imported as n100: primitive rays spanning dimension \(r_l-1\), a unique
strictly positive rational relation, weights proportional to it, an
actual direct lattice and a complete affine factor. The reference and
obstacle classes are its real-PA classes on finite rational spatial
structures, and the same primitive-normal intersection operator and
regularizable domain are used. Thus its real-reference positive theorem
applies when at most one factor is branched.

A bivalent factor has opposite primitive rays in a rank-one lattice,
with equal weights, so is a complete affine line with constant weight.
Merge such factors with \(E\). With one branched factor the imported
one-star-times-affine theorem applies to arbitrary nonseparable
obstacles. With no branched factor the carrier is complete affine space;
n100 includes this case, and includes the point case \(d=0\), where
envelope contact is immediate. No stronger negative quantifier is
inferred from n100. The negative proof above supplies that missing
quantifier directly. This proves the product criterion exactly as
stated.

\subsection{Explicit irrational-reference
illustration}\label{explicit-irrational-reference-illustration}

On two standard unit-weight tripods, \(d=2\), \(C=C'=e=1\), and
\(\kappa=8\). The three standard generators have \(\ell^1\) norms 1, 1
and 2, so the triangle inequality gives \(q\leq2\rho\). Take
\(\theta=\sqrt2/2\) and \[
\gamma_1=\rho,\qquad \gamma_2=2\rho+q+1,\qquad
\gamma=(1-\theta)\gamma_1+\theta\gamma_2
=(1+\theta)\rho+\theta q+\theta.
\] Both components are rational PC positive-growth references, and
\(\gamma\) is genuinely real. Set \(R=1\), \(\lambda=1/4\), \(M_1=3\),
\(M_2=10\), and \[
P=(1-\theta)\max(q/4,\rho-3)
+\theta\max(q/4,2\rho+q-9).
\] On \(\rho\leq2\), both second arguments are at most \(-1\), so
\(P=q/4\). On \(\rho\geq6\), both second arguments dominate, and
\(P=\gamma-(3+7\theta)\). Every displayed maximum has rational walls.
With \(u=P-\gamma+(1-\rho)_+\), the general envelope proof applies and
its full defect is \(-8/16=-1/2\). On the rectangle of the first
standard ray in each star, \(\Delta q=-2\), \(\Delta\rho=0\), so
\(\Delta u=1/2+2\theta\ne0\). This finite example illustrates exact
exterior growth and nonseparability; the universal proof is Lemmas 3--5
and the construction above, not computation on this example.

\subsection{Comparison, limitations, and review
status}\label{comparison-limitations-and-review-status}

The supplied n100 already classifies the product package for rational
references and supplies the positive real-reference side. The prior
real-reference draft gives a negative theorem only under actual ambient
extendibility. The new result fills the remaining universal reference
quantifier on this entire product class. It proves a finite rational
support/decomposition theorem and an admissible exact-growth gluing
construction; it neither assumes nor proves the disputed ambient
Jensen/extension statement. All actual PCCs used in the proof have
carrier centres and incident endpoint cells.

The original P05.2 atom asks for all balanced rational polyhedral pairs.
That remains open. This theorem does not classify finite graph factors
with bounded edges or cycles, arbitrary higher-rank ray stars, general
singular fans, or the proposed intrinsic smoothness/strict-positivity
package. It does not prove all-obstacle regularity for the negative
products. Its conclusion already supplies a regularizable PAPC-envelope
contact defect, so that missing regularity is unnecessary for the
negative statement.

No general approximation theorem for non-PA PC functions or envelopes is
claimed. Lemma 5 decomposes one \textbf{already PA} function on one
fixed finite rational structure, exactly. It supplies no approximation
by rational slopes with bounded individual reference difference. Such an
approximation could change recession behavior, which the exact-sum
gluing avoids.

There is no new literal integral-slope all-obstacle theorem: the real
convex weights and the common rational scale \(\lambda\) do not preserve
integral slopes or rational constants. The previously supplied literal
examples are distinct. The rational-reference criterion remains the
special coefficient scope of n100, while this new theorem genuinely
quantifies over arbitrary real PA PC references in the product class.
There is no analytic realization, CLD/BFJ comparison, nonproper
projection assertion, or P05.5 conclusion. Status is \textbf{partial},
candidate is \textbf{false}. The new proof has its own source-aware
independent PASS; the older exact-scope receipt is not substituted for
that review. The original classification remains open.

\subsection{Reproduction and
bibliography}\label{reproduction-and-bibliography}

The authoritative proof is \texttt{P05.2-argument.md}, with standalone
\texttt{P05.2-partial.tex} and PDF \texttt{build/P05.2-partial.pdf}. The
full author self-audit, research record, exact five-key status, source
manifest and artifact hashes are retained. Run \texttt{bash\ rerun.sh}
in the preserved worker workspace
\path{/workspace/research/campaigns/round-05/P05.2-real-reference-universal-continuation/artifacts/codex-o8EmFf}.
It executes small finite diagnostics, generates
standalone TeX with installed Pandoc, compiles twice with installed
pdfLaTeX and runs artifact checks. Programs, outputs, build commands and
environment details are retained. Compilation checks typesetting;
diagnostics are finite observations, not exhaustive reference testing or
mathematical verification.

No installation, additional model call, collaboration subagent, heavy
computation, network retrieval, external message,
authentication/proxy/CA change or old-workspace write was performed. No
supervised heavy-compute service was supplied or needed. Cached sources
only were used. All old mathematical input bytes are checked unchanged.

\begin{enumerate}
\def\labelenumi{\arabic{enumi}.}
\tightlist
\item
  A. M. Botero, E. Mazzon and L. Pille-Schneider, \emph{Monge--Ampère
  measures on balanced polyhedral spaces}, arXiv:2603.08664v1 (2026).
  Exact definitions and inspected divisor proof passages are stated
  above; needed real-scope arguments are supplied.
\item
  Caller-imported note n100, \emph{P05.2 partial: a branching criterion
  for products of rank-one ray stars}, unchanged supplied argument,
  exact-scope independent PASS receipt. Only its precise real-reference
  positive theorem is imported without reproving it.
\item
  L. Allermann and J. Rau, \emph{First steps in tropical intersection
  theory}, Math. Z. 264 (2010), 633--670, Proposition 3.7. Bibliographic
  explanation of BMPS Lemma 2.35, not an additional unchecked premise of
  the new negative proof.
\item
  R. T. Rockafellar, \emph{Convex Analysis} (1970), Theorems 18.5 and
  19.1. Background for elementary finite generation; Lemma 1 supplies
  the used statements and proofs here.
\item
  The exact dossier in \texttt{input.json} and the previous
  real-reference draft, preserved under \texttt{sources/}. Target
  specifications and fallible research data, not supporting theorems.
\end{enumerate}

\end{document}
